\chapter{The Hydraulic Crawler}
\label{ch:crawler}

A tracked excavator you climb into and drive, with three attachments and a demolition tool on the end
of the arm.

\begin{iefork}
Entirely this fork, and its first vehicle.
\end{iefork}

\begin{ietip}
This is \emph{not} the Excavator in chapter~\ref{ch:multiblocks}. That one is a bucket-wheel
multiblock that mines ore from a mineral vein and does not move. This one is a machine you sit in.
\end{ietip}

\section{Getting one}

Craft the Crawler and place it on level ground with about five blocks clear --- it is a large machine
and it will refuse to go anywhere it does not fit, saying so rather than silently doing nothing.
Right-click to climb in. Sneak and strike it with an Engineer's Hammer to take it away again.

\section{Driving}

\begin{itemize}
\item \textbf{W} and \textbf{S} drive; \textbf{A} and \textbf{D} steer the tracks, which turn on the
      spot the way tracks do.
\item \textbf{Looking around slews the house} --- the cab, the counterweight and the arm all turn
      with your head. You cannot look around independently of the machine, exactly as you cannot in
      the real thing.
\item \textbf{R} and \textbf{F} raise and lower the arm.
\item \textbf{X} and \textbf{Z} reach out and draw back in.
\end{itemize}

The arm is \emph{aimed} rather than jointed: one pair of keys says which way it points, the other how
far along that line the bucket sits, and the boom and stick work out their own angles. Between them
they cover a band from about two blocks out to nearly six, and from five blocks up to nearly three
below the tracks.

\begin{ietip}
The panel in the corner shows the arm's height, its reach, the fuel, and the list of controls.
\end{ietip}

\section{Attachments}

\textbf{G} changes the tool on the end of the arm, \textbf{V} works it.

\begin{itemize}
\item The \textbf{Bucket} digs and keeps what it digs, in nine slots inside the machine. Sneak and
      press \textbf{V} to tip it out; a pipe can also unload it.
\item The \textbf{Grapple} picks up one block --- or one animal --- and puts it down again. A block
      goes back exactly as it came up, facing and all.
\item The \textbf{Breaker} destroys whatever it touches and leaves it on the ground.
\end{itemize}

\begin{iewarning}
The Breaker is a demolition tool and behaves like one. It respects land claims and protected areas
exactly as your own pickaxe does --- every block it takes is on your account --- but within your own
build it will happily take a wall down faster than you can change your mind.
\end{iewarning}

\section{Fuel}

The Crawler runs on \textbf{diesel} and holds eight buckets. Refuel it by hand with any container, or
drive it up to a Gas Station Pump.

Idling costs a little and working costs a lot, so a tank goes a long way if you are only driving.
When the fuel runs low the \emph{attachment} stops but the \emph{tracks keep turning} --- that last
reserve is deliberately kept back so you can always drive home rather than being stranded next to a
half-demolished building.
